% 06-discussion.tex -- Conditional on results that do not exist yet. Text in \pending{} is the planned wording
% and is kept, changed or deleted by the pre-committed rules (protocol.md section 8a).
\section{Discussion}\label{sec:discussion}

\subsection{A Replayed Genuine Trajectory Is In-Distribution}\label{sec:disc-replay}

\todo{R5 decides this subsection. Keep it only if no benign-only detector exceeds 0.60 on DataReplay windows of
$\ge 4$ rows. If all intervals lie below 0.5, keep it and add that inverting the score would need labels. If R4
fires for DataReplay, reduce it to ``not evaluable at this unit''.}
In the DataReplay attack the Sybil pseudonyms broadcast the trajectory of a genuine neighboring vehicle. Seen from
one receiver and one pseudonym, the resulting window is a genuine trajectory under a new pseudonym. A detector that
scores each link on its own compares the window with benign motion, and the window is benign motion: it is inside
the training distribution by construction. \pending{Our results are consistent with this: matched on window
length, no benign-only detector separates DataReplay from benign windows (length-matched AUROC 0.xxx--0.xxx), and
the detectors differ less from each other than from chance.} This is a property of the threat model and of the
per-sender unit, not of one model family. It differs from the cases in which a generative model assigns high
likelihood to unrelated out-of-distribution data~\cite{nalisnick2019deep,kirichenko2020normalizing} or an
autoencoder reconstructs anomalies well~\cite{gong2019memorizing,yoon2021autoencoding}; there the cause is the
model's inductive bias. Replay is also known to be hard for reconstruction-based detection on the in-vehicle
CAN bus~\cite{longari2022candito,hanselmann2020canet}. What can reveal the attack is context across senders, for
example two pseudonyms claiming the same track; no detector in this study has that input.

\subsection{Attack-Majority Training Data}\label{sec:disc-majority}

\todo{R6 decides this subsection. If the mix-gridfolders autoencoder reaches $\ge 0.65$ on GridSybil, delete the
finding from abstract and title and report the number in one sentence. With two confounded regimes, do not write
``the more attack windows, the worse'': that needs E6.}
An unsupervised detector learns what is frequent. Contaminated training data is known to degrade
autoencoders~\cite{zhou2017anomaly,ruff2021unifying}, and methods exist for a small unlabeled anomalous
minority~\cite{qiu2022latent}. In the two GridSybil folders the situation is different in kind: the attack
contributes 64.6\% of the training windows, so the attack is the majority class and ``normal'' is redefined.
\pending{In this regime the convolutional autoencoder scores 0.xxx on GridSybil, against 0.xxx when trained on
the benign windows of the same folders, and the Mahalanobis score falls to 0.xxx.} To our knowledge, a
majority-anomaly training set has not been studied for misbehavior detection. The practical consequence is that
``unsupervised'' on this dataset has two meanings, benign-only data selected with ground-truth labels and
unlabeled recorded traffic, and results obtained under the first do not carry over to the second.

\subsection{Window Length: Signature and Shortcut}\label{sec:disc-length}

\todo{Uses measured values for the control; sentences on detectors follow R4.}
DataReplay and the DoS attacks change pseudonym after one or two messages, so a receiver sees them as many
very short links: 94.9--99.7\% of their windows have at most three rows, and length alone reaches 0.952--0.967.
We regard length in two ways. As a \emph{signature}, a burst of short-lived pseudonyms is real, receiver-observable
evidence of these attacks, and a deployed detector may use it. As a \emph{shortcut}~\cite{geirhos2020shortcut},
it tells nothing about whether a model has learned plausible motion, which is what a sequence model is claimed to
contribute, and it may not hold under another pseudonym-change policy. A pooled AUROC cannot tell the two apart.
Every detector here can see length, so we report each result next to the one-feature control and matched on
length, in line with the call for simple baselines~\cite{arp2022dos,kim2022rigorous}.
\pending{Matched on length, the benign-only detectors retain 0.xxx--0.xxx on the DoS attacks.} The unit of
analysis creates this effect: with pseudonym links, the attack mechanism determines the sequence length. Another
unit, such as all messages a receiver hears in a time interval, would change it and is left to future work.

\subsection{The Gap to Published One-Class Results}\label{sec:disc-gap}

\todo{Observation only. Not allowed without E7 or E9: ``inflated'', ``due to leakage'', any cause. Before
submission: full-text check of the cited papers' metrics, split and sequence unit (gap-analysis.md, Section 6);
the values 0.93--0.99 are from our literature notes.}
Published one-class and reconstruction-based detectors report recall or $F_1$ values of
\pending{0.93--0.99} for Sybil attacks on VeReMi-Extension~\cite{alladi2021deep,kalogiannis2026pampos,%
vishwanath2026federated}. \pending{Under our protocol the best benign-only detector reaches an AUROC of 0.xxx on
GridSybil.} We report this difference as an observation and do not attribute it to a cause. The numbers are not
measured on the same scale (accuracy, recall or $F_1$ at a chosen threshold against threshold-free AUROC), and
the studies differ from ours in at least three respects that we did not vary: the split unit, the source of the
decision threshold, and the sequence unit (sequences grouped by sender identifier against receiver-side
pseudonym links). Any one of them, or the detectors themselves, could account for the difference.

\subsection{Limitations and Threats to Validity}\label{sec:limitations}

\textbf{Near-copies across families (unmeasured).} The GridSybil 0709 folder is a separate re-simulation of the
0709 traffic with different vehicle identifiers, so a vehicle key does not link its vehicles to their counterparts
in the other 0709 folders. Near-copies of some benign test trajectories may therefore be in the training set. We
did not measure this for the split used here; the sensitivity rows without that folder's benign windows and for
the 1416 period alone bound its effect on our numbers.

\textbf{Repeated benign traffic.} Benign traffic appears once per attack family, about four times in total.
Benign windows are thus not independent. The bootstrap resamples vehicle keys, which keeps the copies together,
but the effective benign sample is smaller than the window count suggests.

\textbf{One dataset, one split.} All results are from one simulated dataset and one split seed. The intervals
cover test sampling only; five seeds give a rough standard deviation.

\textbf{Sender-held-out only.} Receivers are not held out: the receiver's position and velocity, four of the 13
features, come from vehicles that also hear training senders.

\textbf{One training attack.} GridSybil is the only attack in any training pool. The supervised results on the
other three attacks are transfer from GridSybil, in one direction, and say nothing general about unseen attacks.

\textbf{No window-level plausibility score.} The rule-based plausibility checks shipped with the
dataset~\cite{kamel2020veremi,kamel2020simulation} decide per message; we have no window-level score for them
under this protocol, so the non-learned reference here is the length control only.

\textbf{Privileged grouping.} In the GridSybil scenarios the ghost vehicles share one pseudonym value; their
messages were separated by the true sender identifier when links were built. A receiver does not have this
information, so GridSybil windows are cleaner than a deployed system would see.

\textbf{Prior exposure and budget.} Two detectors were developed with knowledge of seed-0 test results
(Section~\ref{sec:exposure}), and the LSTM autoencoder has a lower epoch cap than the convolutional one.
\todo{Update after the LSTM epoch-cap decision; report how many runs stopped at the cap.}

\subsection{Out of Scope}\label{sec:scope}

The following were excluded in advance and are not answered here: a sweep over the attack share in the training
data, which would turn the two mix regimes into a curve; a random-split run of the same detectors and a
measurement of the near-copies, which would be needed before any statement about the effect of the split;
leave-one-attack-out training; an ablation of the sequence unit; a window-level plausibility score; and transfer
between the 0709 and 1416 periods. We also did not test learned one-class objectives or contamination-aware
training~\cite{ruff2018deep,ruff2020deep,qiu2022latent}.
